% =====================================================================
%  CHAPTER 2: 儿童少年生长发育概述（对应大纲第一章：生长发育规律）
%  参考PPT：第二章 儿童少年生长发育概述（苏普玉，安徽医科大学）
% =====================================================================
\chapter{儿童少年生长发育概述}
\setcounter{chapter}{2}
\chapterintro{
掌握生长发育基本概念（生长、发育、成熟、成熟度）；掌握生长发育指标体系（体格、体能、心理行为发育指标）；掌握生长发育一般规律（阶段性与连续性统一、程序性与时间性协调、不同步性与多样性平衡、高度可塑性）；掌握头尾发展律和近侧发展律；掌握Scammon生长模式；掌握追赶性生长和生长轨迹现象；熟悉生长发育基本理论观点（进化观、社会生态学理论、整体观）。
}

% ----- 2.1 生长发育概念与指标体系 -----
\section{生长发育概念与指标体系}

\subsection{生长发育基本概念}

\begin{keybox}
\textbf{（一）生长（Growth）}

指身体各部分以及全身在大小、长短和重量上的增加以及身体化学成分的变化。包括：
\begin{itemize}[leftmargin=*]
    \item \imp{形态生长}：细胞、组织、器官在数量、大小和重量上的增加
    \item \imp{化学生长}：细胞、组织、器官、系统的化学成分变化
\end{itemize}
通常使用较多的是涉及形态方面的生长，如身高生长、体重生长、骨骼生长等。

\textbf{（二）发育（Development）}

指身体组织、器官、各系统在功能上的不断分化与完善的过程，包括"身"（体格、体能）、"心"（心理、情绪、行为等）两个密不可分的方面，"发育"在心理学、教育学上也称之为"发展"，社会领域又常用"成长"一词。

\textbf{生长和发育之间的关系：}生长和发育密不可分，生长是发育的前提，发育寓于生长之中。对细胞、组织和器官而言，形态的变化必然伴随功能的分化、增强。因此，常把生长发育一起表述，即身体组织、器官、系统在数量、形态、重量及功能随年龄变化与完善过程。

通常可用"发育"来简称"生长发育"，如"身高生长"在许多情况下用"身高发育"来替代；但一般不能用"生长"代替"发育"，如"视力发育"不能称之为"视力生长"。

\textbf{（三）成熟（Maturity）}

指生长和发育达到一个相对完备的阶段，标志着个体在形态、生理功能、心理素质等方面都已达到成人水平，具备独立生活和养育下一代的能力。

\textbf{成熟度（Maturity Degree）：}专指某一特定生长发育指标当时达到的发育水平占成人水平的百分比，例如，大脑重量在出生时为350$\sim$400g，相当于成人脑重的25\%，而到了6岁，脑重已相当于成人的90\%，但脑发育仍在继续旺盛地进行。

\textbf{生长的终点：}达到成年体型。

\textbf{成熟的终点：}人类在功能上能够成功生育，而不仅仅是能够产生有活力的精子（对于男性）和卵子（对于女性）。因此成功的成熟不仅需要生物学上的成熟，还需要行为和社会层面的成熟。
\end{keybox}

\subsection{生长发育指标体系}

\begin{defbox}
\textbf{（一）体格发育指标}

指身体外部形态的发育，是人体整体发育的重要方面，其测量指标可用人体测量学方法准确测量，大体分为纵向测量指标、横向测量指标、重量测量指标三大类。

\begin{enumerate}[leftmargin=*]
    \item \imp{纵向测量指标}：3岁前有身长、顶臀长、坐高；学龄期有身高、坐高、上肢长、下肢长、手长、足长等。
    \item \imp{横向测量指标}：包括围度和径长。围度主要有头围、胸围、腹围、腰围、上臂围、大腿围和小腿围等；径长主要有肩宽、骨盆宽、胸廓前后径、胸廓左右径、头前后径、头左右径等。
    \item \imp{重量测量指标}：目前在儿童少年卫生/学校卫生工作中主要有体重。
    \item \imp{体格发育的派生指标}：指两项或几项体格发育指标通过数学式表达，用以反映体形和营养状况，如体质量指数（BMI）、Quetelet指数、腰高比和腰臀比等。
\end{enumerate}
\end{defbox}

\begin{keybox}
\textbf{（二）体能发育指标}

体能是健康概念的重要延伸，用以全面、准确评价人体的生理功能和健康状况。学者们将体能细分为健康相关体能和运动相关体能，两者分别以生理功能和运动素质指标反映。

\textbf{1. 健康相关体能指标：}
\begin{itemize}[leftmargin=*]
    \item \imp{心血管功能指标}：一定负荷下人体心率、脉搏、动脉血压的变化
    \item \imp{肺功能指标}：呼吸频率、肺活量、最大通气量、最大吸（摄）氧量等
    \item \imp{肌力发育指标}：握力、背肌力等
\end{itemize}

\textbf{2. 运动相关体能指标：}

指提高运动技能所需要的力量（power）、速度（speed）、灵敏性（agility）、平衡（balance）、协调能力（coordination）和反应时间（reaction time）等方面。

\begin{itemize}[leftmargin=*]
    \item \imp{力量指标}：俯卧撑、引体向上、屈臂悬垂、立定跳远、仰卧起坐、掷铅球、手球掷远、投垒球等
    \item \imp{耐力指标}：20m折返跑（心肺耐力），小学生的50m×8往返跑，大中学男生的1000m跑、大中学女生的800m跑；台阶运动试验（肌耐力）；最大耗氧量测试（全身耐力）
    \item \imp{速度指标}：短跑、球类、游泳、滑雪、击剑、武术等
    \item \imp{灵敏性指标}：10m×4往返跑、反复横跳、蛇形运球、平衡、技巧运动等
    \item \imp{柔韧性指标}：立位体前屈、俯卧位上体上抬等
\end{itemize}

\textbf{3. 体能发育的派生指标：}
\begin{itemize}[leftmargin=*]
    \item \imp{体重肺活量指数}：即肺活量与体重的比值，反映每千克体重的肺活量（ml）
    \item \imp{布兰奇心功指数（Blanche cardiac function index，BI）}：考虑了心率和血压因素，能较全面地反映心脏和血管的功能
\end{itemize}
\end{keybox}

\begin{keybox}
\textbf{（三）心理行为发育指标}

心理行为发育大体分为认知能力、情绪与情感、气质和人格发育，心理行为发育指标通常和心理测验联系在一起。

\textbf{1. 认知能力指标：}
\begin{itemize}[leftmargin=*]
    \item \imp{感知觉}：时间知觉、空间知觉、速度知觉、肌肉用力感、动觉方位感等
    \item \imp{注意力}：集中注意能力、注意广度、注意力分配等
    \item \imp{记忆力}：将学到的信息"储存"和"读出"的神经活动过程，分为感觉记忆、短时记忆、长时记忆等
    \item \imp{思维和想象力}：思维是人脑对客观事物间接、概括的反映，能够认识事物的本质和事物之间的内在联系，属认知的高级阶段。想象是对已有的表象进行加工改造，形成新形象的过程
    \item \imp{语言}：早期词汇生成是指为了交流而说出的单词和形成的手势，词汇理解是指被理解的单词
\end{itemize}

\textbf{2. 情绪与情感指标：}情绪是客观事物是否符合人的需要而产生的态度体验，反映客观事物与人的需要之间的关系。情感则是与人的高级社会性需要满足与否相联系的态度体验，是在社会交往的实践中逐渐形成的，如友谊感、道德感、美感和理智感等。

\textbf{3. 气质与人格：}气质是心理过程的强度、速度、稳定性与指向性等方面在行为上的表现。人格又称个性，是个体在适应社会的成长过程中，经遗传与环境的交互作用形成的稳定而独特的心理特征，气质和性格是人格最重要的部分。
\end{keybox}

\subsection{生长发育研究内容}

\begin{defbox}
\textbf{（一）不同层次水平上的身体发育研究：}
\begin{enumerate}[leftmargin=*]
    \item \imp{整体水平}：体格、体形、体姿、生理功能、运动素质等发育随年龄发生的变化
    \item \imp{器官-系统水平}：研究呼吸系统、听觉器官、性器官等形态变化和功能成熟过程及其影响因素
    \item \imp{体成分水平}：研究体水分、肌肉质量增加、脂肪沉积等性别二态性差异表现
    \item \imp{组织学水平}：重点关注上皮组织、结缔组织、肌肉组织和神经组织的形态结构和功能变化过程
    \item \imp{细胞水平}：研究细胞生长、细胞分裂与分化、细胞更新等生命过程
    \item \imp{分子水平}：从遗传物质量、基因变异、基因非变异的可遗传性表达（表观遗传）等水平
\end{enumerate}

\textbf{（二）心理行为发育：}认知和情绪发育、个性特征发育、社会行为发育。

\textbf{（三）青春期发育：}
\begin{itemize}[leftmargin=*]
    \item 青春期启动的机制尚未揭示：目前已知下丘脑GnRH脉冲分泌引发的垂体促性腺激素分泌模式的变化，是青春期启动的触发因素
    \item 生长发育长期变化（生长发育长期趋势）是指以儿童体格发育、青春期和成年身高为核心表现的长期变化
\end{itemize}
\end{defbox}

% ----- 2.2 生长发育一般规律 -----
\section{生长发育一般规律}

\subsection{生长发育的阶段性与连续性统一}

\begin{keybox}
\textbf{（一）生长发育的阶段性}

生长发育不只是量的增加，还有质的变化，因而形成不同的发育阶段。前阶段为后一阶段的发育准备了基础，而其出现的发育障碍，必然对后阶段产生不良影响。生长发育的各种表现存在关键期或敏感期的现象，在不同阶段，其发展任务应有所侧重。

\textbf{发育任务（Developmental Task）：}由美国生理学家和教育心理学家罗伯特·詹姆斯·哈维格斯特（Robert J. Havighurst）1972年提出，在一定的年龄段，个体的心理行为成熟程度应能达到一定的水平。这些发育任务是特定年龄的基本任务，既是教养目标，也是判断发育水平的依据。

\textbf{（二）生长发育的连续性}

生长发育是一个动态的连续过程，这一过程是量的积累和功能成熟，按照遗传的潜能在生长发育。

\textbf{1. 生长的轨迹现象（Canalization/Growth Trajectory）：}

生长轨迹或称之为生长管道现象是人体生长发育过程的一种生物学现象，当外界环境无特殊变化的条件下，个体生长发育过程较为稳定，呈现出轨迹现象，使生长中的个体在群体中保持有限的上下波动幅度。个体儿童的生长轨迹既与遗传有关，还与疾病及其治疗、营养、体力活动、情感环境等多种因素有关。

\textbf{2. 追赶性生长（Catch-up Growth）：}

当处在生长发育过程中的个体受到疾病、营养、心理应激等因素作用下，会出现暂时性生长发育的连续性被打破现象，如生长迟缓，一旦这些影响因素解除，机体表现为向原有的正常轨迹靠近并具有生长发育强烈的倾向，年龄越小越明显，生长的加速幅度越大。这种阻碍儿童生长的病理性因素被克服或解除后表现出的加速生长并逐渐恢复到正常轨迹的现象称为追赶性生长。
\end{keybox}

\subsection{生长发育的程序性和时间性协调}

\begin{keybox}
\textbf{（一）生长发育的程序性}

\textbf{1. 运动能力发育的程序性：}

\begin{itemize}[leftmargin=*]
    \item \imp{头尾发展律（Cephalocaudal Development）}：胎儿至婴幼儿期体格和粗大动作发育遵循从头部到尾部发展的程序性。粗大动作：按抬头、翻身、坐、爬、站、走、跑、跳的发育程序进行。
    \item \imp{近侧发展律（Proximodistal Development）}：胎儿期至儿童期体格发育还遵循从身体的近端（或中心）到远端（或外周）发展的程序性，表现为躯干的生长先于四肢，四肢的近端生长先于远端。
\end{itemize}

\textbf{2. 线性生长的规律：}

儿童出生后有婴儿期、童年期和青春发动期三个生长模式，婴儿期生长可能延续了胎儿期生长突增，推测婴儿期生长主要与胎儿营养有关；童年期生长一般开始于出生后6$\sim$12个月，典型表现为生长的突然加速，可能与生长激素开始起作用有关。

\textbf{（二）生长发育的个体差异性}

上述的生长发育指标呈现的程序特征，反映了人类在生长发育过程所表现的共同特征，当环境因素发挥作用或遗传和环境共同作用时，儿童少年身体和心理发育方面的表现则存在明显的个体差异。
\end{keybox}

\subsection{生长发育的不同步性与多样性平衡}

\begin{keybox}
\textbf{（一）不同组织器官系统的生长不同步性——Scammon生长模式}

不同的组织分化发育成器官，不同的器官相互结合完成特定的功能便构成不同的系统。不同器官、系统的生长发育的类型称为生长模式。理查德·斯卡蒙（Richard Scammon）通过发育水平曲线的描述，发现人体器官、系统的生长曲线可分为四种类型：

\begin{table}[H]
\centering
\caption{Scammon生长模式的四种类型}
\begin{tabularx}{\textwidth}{lX}
\hline
\textbf{类型} & \textbf{特点} \\
\hline
\imp{一般型（淋巴系统除外的器官）} & 婴幼儿期发育快；学龄前期和学龄期发育平缓；青春期再次加快；青春期后逐渐停止 \\
\imp{淋巴系统型} & 淋巴结、扁桃体、胸腺等淋巴器官——儿童期发育很快，青春期达高峰，此后逐渐衰退，成人时仅相当于高峰期的一半 \\
\imp{神经系统型} & 中枢神经系统、脊髓、眼球等——出生后有一个生长突增期，婴幼儿期和学龄前期发育迅速，6岁时脑重已达成人90\%左右 \\
\imp{生殖系统型} & 睾丸、卵巢、子宫等生殖器官——青春期前基本处于停滞状态，青春期开始迅速加快发育直至成熟 \\
\hline
\end{tabularx}
\end{table}

\textbf{子宫型生长模式（附加模式）：}子宫、肾上腺等器官不同于上述生长模式，其在出生时较大，其后迅速变小，青春期开始前才恢复到出生时的大小；其后迅速增大。

\textbf{（二）生长发育多样性}

生长发育是受生物遗传、环境污染和气候变化、心理和社会背景等多种因素的共同作用。生长发育的路径千差万别，每种生长发育现象有多种潜在的路径和结果。生长发育的多样性可表现为体格、体能、认知发育的个体差异。
\end{keybox}

\subsection{生长发育的高度可塑性}

\begin{defbox}
\textbf{发育可塑性（Developmental Plasticity）：}指生物体在生命周期中受外部环境因素或其他内部生理因素影响，通过基因-环境相互作用而发生的发育轨迹的永久性行为、解剖或生理变化，为生物体在不断变化的环境下提供生存和繁殖成功的最佳机会。
\end{defbox}

% ----- 2.3 生长发育基本理论观点 -----
\section{生长发育基本理论观点}

\subsection{生长发育进化观}

\begin{defbox}
\textbf{1. 进化形成的适应机制在生长发育中的作用：}查尔斯·罗伯特·达尔文等进化论学者都特别强调人对进化形成的适应机制的重要性，也强调促使该机制得以激活、发展的环境因素的重要作用。

\textbf{2. 生命史策略在生长发育中的作用：}生命史理论研究生物体在不同的生命目标（如生存、生长和繁殖）之间进行权衡并适应性地分配资源的理论。个体在资源有限的情况下，需要考虑如何分配自身资源，而个体早期成长环境会通过内分泌轴，影响资源在生存、生长和繁衍的分配（即生命史策略）。
\end{defbox}

\subsection{社会生态学理论}

\begin{keybox}
\textbf{Bronfenbrenner生态系统理论：}

生态系统理论是发展心理学的重要理论之一，由美国心理学家尤里·布朗芬布伦纳（Urie Bronfenbrenner，1917—2005）建立的个体发展模型。他认为个体的发展嵌套于相互影响的一系列的环境系统中，各个系统之间存在交互耦合的联系和影响，个体与环境相互作用并影响着个体的发展。

\textbf{（一）空间维度：}
\begin{enumerate}[leftmargin=*]
    \item \imp{微系统（Microsystem）}：是个体活动和人际交往的直接环境，不断在变化和发展，包括家庭、学校、同伴群体、社区等。
    \item \imp{中系统（Mesosystem）}：指各微系统间的联系或相互关系，如家庭、同伴和学校之间的关系。
    \item \imp{外系统（Exosystem）}：指儿童并未直接参与但却对他们的发展产生影响的环境。
    \item \imp{宏系统（Macrosystem）}：指存在于以上三个系统中的文化、亚文化和社会环境。
\end{enumerate}

\textbf{（二）时间维度（历时系统，Chronosystem）：}

布朗芬布伦纳还把时间维度纳入模型，将时间作为研究个体成长中的心理变化参照体系，称为历时系统。他强调儿童的发展、变化是一个将时间和环境结合起来考察的动态过程。生态系统不是一种恒定不变的系统，它始终处于动态变化中。
\end{keybox}

\subsection{整体观}

\begin{defbox}
\textbf{生长发育整体观}主张将儿童少年的生长发育人为地划分成几个学科领域，不同学科领域在揭示生长发育现象时，本质相同但各有侧重点。生长发育整体观理论的核心观念是：应将儿童少年生长发育的各个方面都视为是存在有机联系的，相互依赖的，而不是零碎的、割裂的。每个人在生物性、认知性和社会性等方面都具有自身的特征，每个人又受制于各个层次系统"大环境"的直接间接的影响。
\end{defbox}

% ----- 复习题 -----
\section{复习题}

\begin{enumerate}[leftmargin=*, itemsep=8pt]
    \item \textbf{【名词解释】}生长；发育；成熟；成熟度
    \item \textbf{【名词解释】}头尾发展律；近侧发展律；追赶性生长；生长轨迹现象
    \item \textbf{【名词解释】}Scammon生长模式；发育可塑性
    \item \textbf{【名词解释】}发育任务；生态系统理论（Bronfenbrenner）
    \item \textbf{【简答题】}简述生长与发育的区别与联系。
    \item \textbf{【简答题】}简述生长发育指标体系及其主要内容。
    \item \textbf{【简答题】}试述头尾发展律和近侧发展律的具体表现。
    \item \textbf{【简答题】}简述Scammon生长模式中四种类型的特点。
    \item \textbf{【简答题】}何谓追赶性生长？试从正反两个方面阐述其意义。
    \item \textbf{【简答题】}简述Bronfenbrenner生态系统理论的空间维度和时间维度。
    \item \textbf{【简答题】}试述生长发育生态观中"生态系统理论——空间维度"在理解儿童少年网络成瘾行为中的意义。
    \item \textbf{【简答题】}简述生长发育整体观的核心观念。
\end{enumerate}

\subsection*{参考答案要点}

\begin{enumerate}[leftmargin=*, itemsep=5pt]
    \item \textbf{生长}：身体各部分以及全身在大小、长短和重量上的增加以及身体化学成分的变化，包括形态生长和化学生长。\textbf{发育}：身体组织、器官、各系统在功能上的不断分化与完善的过程，包括"身"（体格、体能）、"心"（心理、情绪、行为等）。\textbf{成熟}：生长和发育达到一个相对完备的阶段，标志个体在形态、生理功能、心理素质等方面达到成人水平。生长和发育密不可分，生长是发育的前提，发育寓于生长之中。

    \item \textbf{头尾发展律}：胎儿至婴幼儿期体格和粗大动作发育遵循从头部到尾部发展的程序性——抬头→翻身→坐→爬→站→走→跑→跳。\textbf{近侧发展律}：从身体近端（中心）到远端（外周）发展——躯干生长先于四肢，四肢近端先于远端。\textbf{追赶性生长}：阻碍生长的病理性因素被克服后表现出的加速生长并逐渐恢复到正常轨迹的现象，年龄越小越明显。\textbf{生长轨迹现象}：外界环境无特殊变化时，个体生长发育过程较为稳定，呈现出轨迹现象。

    \item \textbf{Scammon生长模式}：四种类型——一般型（婴幼儿期快→学龄期平缓→青春期加速）、淋巴系统型（儿童期快→青春期高峰→成人期衰退）、神经系统型（婴幼儿期迅速→6岁接近成人90\%）、生殖系统型（青春期前停滞→青春期迅速）。另有子宫型生长模式（出生时大→迅速缩小→青春期恢复并增大）。\textbf{发育可塑性}：生物体受外部环境或内部生理因素影响，通过基因-环境相互作用发生的发育轨迹的永久性变化。

    \item \textbf{发育任务}：由Havighurst 1972年提出，在一定的年龄段，个体心理行为成熟程度应能达到一定水平，既是教养目标也是判断发育水平的依据。\textbf{生态系统理论}：Bronfenbrenner建立的个体发展模型，认为个体发展嵌套于相互影响的一系列环境系统中（微系统→中系统→外系统→宏系统），并包含时间维度（历时系统）。

    \item 见第1题。生长是形态和化学变化，发育是功能分化和完善。生长是发育的前提，发育寓于生长之中。可用"发育"简称"生长发育"，但不能用"生长"代替"发育"。

    \item 三大类：(1)体格发育指标（纵向：身高、坐高等；横向：围度、径长；重量：体重；派生指标：BMI等）；(2)体能发育指标（健康相关：心血管、肺功能、肌力；运动相关：力量、耐力、速度、灵敏、柔韧；派生指标：体重肺活量指数、BI等）；(3)心理行为发育指标（认知能力、情绪与情感、气质与人格）。

    \item 头尾发展律：粗大运动从头到脚——抬头→翻身→坐→爬→站→走→跑→跳。近侧发展律：从躯干近端向远端发展——躯干→上肢近端→上肢远端→手指精细动作。

    \item 见第3题Scammon生长模式。

    \item 追赶性生长定义见第2题。正面意义：说明生长发育有强大的自我调节和恢复能力，早期干预对恢复生长发育至关重要。反面意义：追赶性生长过度可能导致代谢编程改变，增加肥胖和代谢综合征风险。

    \item 空间维度四层次：微系统（家庭、学校、同伴等直接环境）→中系统（微系统间的联系）→外系统（未直接参与但产生影响的环境）→宏系统（文化、亚文化和社会环境）。时间维度（历时系统）：将时间作为参照体系，强调儿童发展变化是将时间和环境结合考察的动态过程。

    \item 从微系统看：同伴群体中的网络使用规范和压力。从中系统看：家庭与学校在网络使用管理上的协调。从外系统看：父母工作压力导致减少监管。从宏系统看：数字文化和社会规范对网络行为的塑造。

    \item 应将儿童少年生长发育的各个方面都视为存在有机联系、相互依赖的，而不是零碎的、割裂的。每个人在生物性、认知性和社会性等方面都具有自身特征，又受制于各层次系统"大环境"的直接间接影响。
\end{enumerate}

\newpage
